% talk about fusion
% mention Effuse 
Arising from the concept of ensemble learning in the 1990s, recent works have explored ensemble techniques for Speech Foundation Models (SFMs). \cite{arunkumar22b_SSLFusion} was an early work studying the effect of fusion for self-supervised speech models. 
However, their work only considers performing fusion on the last layer output of each model.
\cite{tang22_fusion} further improves it by extracting multiple layer features from each model while doing fusing.
% Building upon this,~\cite{tang22_fusion} extracts features from multiple layers of each model and discussed different level of fusion. \david{I don't understand what the previous sentence is saying - Ian, can you rephrase it?}
In~\cite{fu23b_otf}, instead of fusing the layer outputs of multiple SFMs, the authors propose to use Optimal Transport to perform the fusion of model weights directly.
While this method significantly reduces computational cost at inference time, it requires extensive end-to-end fine-tuning on the merged model to achieve competitive performance.
Due to computational constraints, we do not include this approach in our study.
Effuse~\cite{srivastava24_effuse}, on the other hand, trains an auxiliary prediction head to use features from one SFM to predict the output of another SFM, eliminating the need to run inference on both models at testing time.
However, their fine-tuning process consists of two stages: first training the upstream models separately, followed by training the predictor.
Due to this multi-stage fine-tuning requirement, we exclude this method from our experiments.
Nevertheless, our baseline method of using weighted sum can be considered to be equivalent to the first stage of Effuse.
Finally, many of these prior works focus exclusively on ASR tasks, while our work also evaluates speaker-related and paralinguistic tasks.

Several papers have also studied layer fusion for SSL speech models. In~\cite{huo23_hff}, the authors proposed methods to fuse layers in a hierarchical way and~\cite{huo23b_gaff} propose a global attentional feature fusion.
% Nonetheless, despite the fact that they allow interaction between the layers, the core of their idea is projection after their fusion module.
More recently,~\cite{chiu24_sinica,shih24_interface} systematically investigated multiple layer fusion strategies for utilizing Speech SSL models.
The former focused on exploring various combinations of layer and model fusion, whereas the latter concentrated specifically on layer fusion.
In terms of evaluation, the former evaluates only on LibriSpeech ASR, while the latter evaluates on a variety of tasks from SUPERB and ML-SUPERB.
% Notice that there are some prior works that also study the similar concept of layer fusion.

% talk about layer aggregation